\documentclass[11pt]{article}
\usepackage[margin=1in]{geometry}
\usepackage{amsmath,amssymb,amsthm}
\usepackage{booktabs}
\usepackage{hyperref}
\newtheorem{theorem}{Theorem}
\newtheorem{lemma}[theorem]{Lemma}
\newtheorem{proposition}[theorem]{Proposition}
\theoremstyle{definition}
\newtheorem{definition}[theorem]{Definition}
\theoremstyle{remark}
\newtheorem{remark}[theorem]{Remark}

\title{Disproof of Conjecture 00000003952:\\
twin-width at most 1 is not the same as distance-hereditary}
\author{jilint777}
\date{2026-10-04}

\begin{document}
\sloppy
\maketitle

\begin{abstract}
Conjecture 00000003952 asserts, as its first clause, that the graphs of twin-width at
most $1$ are exactly the distance-hereditary graphs. Both inclusions are false. The
house (the complement of the path $P_5$) has twin-width $1$ but is not
distance-hereditary. The net (a triangle with a pendant vertex at each corner) is
distance-hereditary but has twin-width $2$. Since the conjecture is a conjunction, it is
false. All definitions (trigraphs, contractions, contraction sequences, twin-width,
walks in induced subgraphs, distance-hereditary graphs) are formalized from scratch in
Lean~4 (core library only). Both counterexamples are proved there, and
\texttt{verify.py} checks them independently.
\end{abstract}

\section{The conjecture and the clause we refute}
\begin{quote}
\emph{Definition:} The twin-width is the minimal width over contraction sequences of
red-black graphs. \emph{Conjecture:} The graphs of twin-width at most 1 are exactly the
distance-hereditary graphs; recognizing twin-width at most 2 is polynomial time, while
recognizing twin-width at most 3 is NP-hard, so the boundary lies exactly at 2.
\end{quote}
The Chinese version states the same. The conjecture is a conjunction of three
statements. We refute the first one,
\[
  \textbf{(C)}\qquad \{G : \operatorname{tww}(G)\le 1\} \;=\; \{G : G \text{ is distance-hereditary}\},
\]
and this already makes the whole conjecture false. We refute \emph{both} inclusions of (C),
each with an explicit connected graph. We do not address the two complexity statements.

\section{Definitions}
All graphs are finite and simple, and have vertex set $\{0,\dots,n-1\}$.

\begin{definition}[Trigraphs and contractions; Bonnet--Kim--Thomass\'e--Watrigant]
A \emph{trigraph} (red--black graph) is a vertex set with, for each pair of distinct
vertices, one of three colours: \emph{none}, \emph{black} or \emph{red}. A graph is the
trigraph with black edges and no red edges. \emph{Contracting} two distinct vertices
$u,v$ replaces them by a single new vertex $z$. Pairs not involving $z$ keep their
colours. For every other vertex $w$, the pair $zw$ is
\begin{itemize}
\item black if $uw$ and $vw$ are both black,
\item none if $uw$ and $vw$ are both none,
\item red otherwise (in particular whenever $uw$ or $vw$ is red).
\end{itemize}
The \emph{red degree} of a vertex is its number of red edges. A \emph{contraction
sequence} of an $n$-vertex graph $G$ is a sequence $G=T_n,T_{n-1},\dots,T_1$ of
trigraphs, where $T_{i-1}$ is obtained from $T_i$ by contracting two vertices. Its
\emph{width} is the maximum red degree of a vertex over all $T_i$. The \emph{twin-width}
$\operatorname{tww}(G)$ is the minimum width over all contraction sequences of $G$. So
$\operatorname{tww}(G)\le d$ iff some contraction sequence has width at most $d$.
\end{definition}

Equivalently, one can follow the partition of $V(G)$ into the sets of original vertices
represented by the current vertices. Two parts $P\ne Q$ are joined by a red edge iff the pair
$(P,Q)$ is not homogeneous (some but not all pairs in $P\times Q$ are edges). By induction,
this gives exactly the trigraphs above, because the merging rule is the same.
\texttt{verify.py} uses this partition formulation and Lean uses the trigraph one.

\begin{definition}[Distance-hereditary; Howorka]
For $S\subseteq V(G)$, $G[S]$ is the induced subgraph and $d_{G[S]}$ its distance
($\infty$ between different components). $G$ is \emph{distance-hereditary} (DH) if
$d_{G[S]}(x,y)=d_G(x,y)$ for every $S$ with $G[S]$ connected and all $x,y\in S$, that is,
if every connected induced subgraph is isometric.
\end{definition}
By Bandelt and Mulder, a connected graph is DH iff it has no induced house, hole $C_k$
($k\ge5$), domino or gem. It is also DH iff it can be reduced to a single vertex by
repeatedly deleting a pendant vertex or one vertex of a pair of twins. We prove everything
from Howorka's definition. The other two characterizations are used only as cross-checks in
\texttt{verify.py}.

\section{The house: twin-width 1, not distance-hereditary}
Let $H$ be the graph on $\{0,1,2,3,4\}$ with edges
\[ 02,\ 03,\ 04,\ 13,\ 14,\ 24 . \]
This is the house: the 4-cycle $0\,3\,1\,4\,0$ with the roof vertex $2$ adjacent to
$0$ and $4$. Its non-edges $01,12,23,34$ form the path $0\,1\,2\,3\,4$, so
$H=\overline{P_5}$.

\begin{theorem}\label{thm:house}
$\operatorname{tww}(H)=1$, and $H$ is not distance-hereditary.
\end{theorem}

\begin{proof}
\emph{$\operatorname{tww}(H)\le1$.} We contract $\{3,4\}$, then add $2$, then $1$, then $0$.
Here $N(0)=\{2,3,4\}$, $N(1)=\{3,4\}$, $N(2)=\{0,4\}$, $N(3)=\{0,1\}$ and
$N(4)=\{0,1,2\}$.
\begin{center}
\begin{tabular}{llll}
\toprule
step & vertices & red edges & max red degree\\
\midrule
$T_5=H$ & $0,1,2,3,4$ & none & 0\\
contract $3,4\to A$ & $0,1,2,A$ & $A2$ & 1\\
contract $A,2\to B$ & $0,1,B$ & $B1$ & 1\\
contract $B,1\to C$ & $0,C$ & $C0$ & 1\\
contract $C,0$ & one vertex & none & 0\\
\bottomrule
\end{tabular}
\end{center}
Details:
\begin{itemize}
\item $A$: $3,4$ are both adjacent to $0$ and to $1$ (black), and only $4$ is adjacent to
$2$ (red).
\item $B$: $A0$ and $20$ are black, so $B0$ is black. $A1$ is black and $21$ is none,
so $B1$ is red. The red edge $A2$ is now inside $B$.
\item $C$: $B0$ is black and $10$ is none, so $C0$ is red.
\end{itemize}
Every trigraph has maximum red degree at most $1$.

\emph{$\operatorname{tww}(H)\ge1$.} Contracting $u,v$ gives the new vertex a red edge to
every $w\notin\{u,v\}$ adjacent to exactly one of $u,v$. If there is no such $w$, then
$u,v$ are twins. $H$ has no twins: the sets $N(x)\setminus\{y\}$ and
$N(y)\setminus\{x\}$ above differ for all $x\ne y$, for instance
$N(0)\setminus\{4\}=\{2,3\}\ne\{1,2\}=N(4)\setminus\{0\}$. So every first contraction
creates a red edge, and every contraction sequence has width $\ge1$.

\emph{Not DH.} Let $S=\{1,2,3,4\}$. Then $H[S]$ has edges $13,14,24$, so it is the
path $2\,4\,1\,3$ and is connected. Moreover $d_H(2,3)=2$ via $2\,0\,3$, since $2$ and $3$
are not adjacent. In $H[S]$ the only path from $2$ to $3$ is $2\,4\,1\,3$, so
$d_{H[S]}(2,3)=3\ne2$.
\end{proof}

\section{The net: distance-hereditary, twin-width 2}
Let $N$ be the graph on $\{0,\dots,5\}$ with edges
\[ 01,\ 12,\ 02,\ 03,\ 14,\ 25 . \]
This is the triangle $0,1,2$ with pendant vertices $3,4,5$ attached to $0,1,2$.

\begin{theorem}\label{thm:net}
$N$ is distance-hereditary, and $\operatorname{tww}(N)=2$.
\end{theorem}

\begin{proof}
\emph{DH.} An induced path of $N$ uses at most one triangle edge: if it used two, say
$a\,b\,c$, then $ac$ would be a chord. A pendant vertex has degree $1$, so it can only be
an endpoint of a path. So an induced path from $x$ to $y$ consists of, in order:
\begin{itemize}
\item the edge from $x$ to its triangle vertex, if $x$ is a pendant;
\item the triangle edge between the two triangle vertices, if they differ;
\item the edge to $y$, if $y$ is a pendant.
\end{itemize}
Hence all induced $x$--$y$ paths of $N$ have the same length. A shortest path is
induced, so this common length is $d_N(x,y)$. Now let $N[S]$ be connected and
$x,y\in S$. A shortest $x$--$y$ path in $N[S]$ is induced in $N[S]$, hence induced in $N$, so its length is $d_N(x,y)$. Therefore
$d_{N[S]}(x,y)=d_N(x,y)$.
(Alternatively: delete the pendants $3,4,5$; the remaining triangle reduces by twin
deletions.)

\emph{$\operatorname{tww}(N)\ge2$.} By the remark in the proof of
Theorem~\ref{thm:house}, the new vertex gets a red edge to every vertex of
$\Delta(u,v)=(N(u)\setminus\{v\})\,\triangle\,(N(v)\setminus\{u\})$. Up to the symmetry of
$N$, which permutes the three corners together with their pendants, there are four types of
pairs:
\[
\Delta(0,1)=\{3,4\},\quad \Delta(0,3)=\{1,2\},\quad \Delta(0,4)=\{2,3\},\quad \Delta(3,4)=\{0,1\}.
\]
These are triangle--triangle, corner with its own pendant, corner with another pendant,
and pendant--pendant. In every case the first contraction already creates a vertex of
red degree $2$. So no contraction sequence has width $\le1$.

\emph{$\operatorname{tww}(N)\le2$.} Contract $3$ into $0$, $4$ into $1$, $5$ into $2$,
then $1$ into $0$, then $2$ into $0$. The maximum red degrees along this sequence are
$0,2,2,2,1,0$.
\end{proof}

\begin{theorem}\label{thm:main}
Clause (C) is false in both directions: $H$ has $\operatorname{tww}\le1$ but is not DH, and
$N$ is DH but has $\operatorname{tww}>1$. Hence Conjecture 00000003952 is false.
\end{theorem}

\section{Interpretation and robustness}
\begin{itemize}
\item \textbf{``Exactly''} means equality of the two graph classes. Each inclusion fails
separately (Theorems~\ref{thm:house} and~\ref{thm:net}).
\item \textbf{Twin-width.} We use the standard definition. The trigraph and partition
formulations agree, as explained above. ``Red-black graphs'' are the trigraphs, and a
graph is a trigraph without red edges. The house sequence has red degree $\le1$ at
\emph{every} vertex of every trigraph. The net lower bound comes from the red degree of
the \emph{newly created} vertex in the first trigraph. So both results also hold for
variants that measure only the contracted vertex, or all vertices. Twin-width does not
depend on vertex names, and in Lean the merged vertex keeps the name of the first one.
\item \textbf{Distance-hereditary.} Both graphs are connected, so the usual definitions
coincide on them: Howorka's isometric-subgraph definition, the Bandelt--Mulder
forbidden-subgraph and pruning characterizations, and the condition that every induced
path is a shortest path. The house is itself one of the forbidden induced subgraphs.
\item \textbf{Context.} Twin-width is invariant under complementation, but the DH
property is not: $\overline{H}=P_5$ is DH. This explains why (C) cannot hold.
\verb|verify.py| also enumerates all connected graphs on at most $6$ vertices:
\begin{itemize}
\item on $5$ vertices, exactly two have $\operatorname{tww}\le1$ but are not DH (the
house and the gem);
\item on $6$ vertices there are $27$ such graphs, and the net is the only DH graph with
$\operatorname{tww}>1$.
\end{itemize}
\item \textbf{Not addressed.} The complexity statements about recognizing twin-width at
most $2$ and at most $3$ are not needed and not discussed.
\end{itemize}

\section{Lean formalization}
The directory \texttt{lean4/} is a self-contained Lean~4.19.0 project (core library only;
no \texttt{sorry}, no \texttt{native\_decide}, no added axioms; \verb|#print axioms|
reports at most \texttt{propext} and \texttt{Quot.sound}). Contents of \texttt{Main.lean}
(namespace \texttt{TwinWidth}):
\begin{itemize}
\item \texttt{Graph} (vertex count and edge list), \texttt{Graph.adj}, \texttt{Graph.WF}
(well-formed simple graph).
\item \texttt{Col}, \texttt{Trigraph}, \texttt{Trigraph.contract}: the contraction rule
above, via \texttt{mergeCol}. Also \texttt{redDeg}, \texttt{maxRedDeg},
\texttt{Graph.toTrigraph} (all edges black).
\item \texttt{IsContractionSeq}: every step contracts two distinct current vertices, and
at most one vertex remains at the end.
\item \texttt{trigraphsAlong} and \texttt{seqWidth}: the maximum red degree over all
trigraphs of the sequence.
\item \texttt{TwwLe G d} $:=\exists s$, $s$ is a contraction sequence of $G$ with
\texttt{seqWidth} $\le d$. \texttt{TwwEq G d} means \texttt{TwwLe G d} and not
\texttt{TwwLe G d'} for $d'<d$.
\item \texttt{Reach G m k x y}: an inductive walk predicate. It says there is a walk
$x=v_0,\dots,v_j=y$ with $j\le k$ and $v_1,\dots,v_j$ in the vertex set encoded by the
bit mask $m$, i.e.\ $d_{G[m]}(x,y)\le k$.
\item \texttt{ConnectedOn G m}, and \texttt{DH G}: for all masks $m<2^n$ with $G[m]$
connected, all $x,y\in m$ and all $k$, \texttt{Reach G (full G) k x y} implies
\texttt{Reach G m k x y}.
\item Computation:
\begin{itemize}
\item \texttt{reach\_iff\_ball} proves that \texttt{Reach} agrees with a computable BFS
ball.
\item \texttt{dh\_of\_check} proves, for every well-formed graph, that a finite Boolean
check \texttt{dhCheck} implies \texttt{DH}. The check includes stabilization of the balls
at radius $n$.
\item \texttt{not\_twwLe\_of\_first} proves that if every first contraction has red degree
$>d$, then \texttt{TwwLe G d} fails.
\end{itemize}
\item Results:
\begin{itemize}
\item \texttt{house\_twwLe\_one}, \texttt{house\_seq\_widths} (red degrees
$[0,1,1,1,0]$), \texttt{house\_tww\_eq\_one}.
\item \texttt{house\_not\_DH}, with witness $m=\{1,2,3,4\}$, $x=2$, $y=3$, $k=2$.
\item \texttt{net\_DH}, \texttt{net\_not\_twwLe\_one}, \texttt{net\_tww\_eq\_two}.
\item Sanity checks: \texttt{P5\_twwLe\_one} and \texttt{P5\_DH}. These show that both
predicates are satisfiable and agree on $P_5$.
\end{itemize}
\item Final theorems:
\begin{itemize}
\item \texttt{twwToDH\_false : \textasciitilde ClauseTwwToDH}
\item \texttt{DHToTww\_false : \textasciitilde ClauseDHToTww}
\item \texttt{conjecture\_00000003952\_false : \textasciitilde Clause}, where\\
\texttt{Clause := $\forall$ G, G.WF $\to$ (TwwLe G 1 $\leftrightarrow$ DH G)}
\item \texttt{conjecture\_00000003952\_conjunction\_false (Rest) : \textasciitilde (Clause
$\wedge$ Rest)}
\end{itemize}
\end{itemize}
Lean does not cover the enumeration of all graphs on at most $6$ vertices or the
equivalence with other DH characterizations. These appear only in \texttt{verify.py} and
are not needed for the disproof.

\section{Independent check and reproduction}
\texttt{verify.py} uses only the Python standard library (plus \texttt{nauty-geng} for the
optional enumeration, if installed). It checks:
\begin{enumerate}
\item the explicit sequences, by trigraph simulation;
\item the exact twin-width of $H$, $N$ and $P_5$, by exhaustive search over partition
sequences;
\item the DH property, in three ways: BFS distances in all induced subgraphs,
Bandelt--Mulder pruning, and forbidden induced subgraphs;
\item the enumeration mentioned above.
\end{enumerate}
\begin{verbatim}
cd lean4 && lake build     # Lean 4.19.0, core library only
cd .. && python3 verify.py
\end{verbatim}

\begin{thebibliography}{9}
\bibitem{bktw} \'E.~Bonnet, E.~J.~Kim, S.~Thomass\'e, R.~Watrigant,
\emph{Twin-width I: tractable FO model checking}, J.~ACM 69 (2022).
\bibitem{howorka} E.~Howorka, \emph{A characterization of distance-hereditary graphs},
Quart.\ J.\ Math.\ Oxford 28 (1977) 417--420.
\bibitem{bm} H.-J.~Bandelt, H.~M.~Mulder, \emph{Distance-hereditary graphs},
J.~Combin.\ Theory Ser.~B 41 (1986) 182--208.
\end{thebibliography}
\end{document}
